\chapter{常用记号}
\label{app:notation}

本附录汇集正文常用的记号；各题中的定义域和约定以题设为准。

\section{数集与运算记号}
\(\mathbb N,\mathbb N^*,\mathbb Z,\mathbb Q,\mathbb R,\mathbb C\) 分别表示非负整数集、正整数集、整数集、有理数集、实数集和复数集。有限集 \(A\) 的元素个数记为 \(|A|\)。集合运算及逻辑记号见附录\ref{app:sets-logic}。

\(\sum_{k=1}^{n}a_k=a_1+\cdots+a_n\)，\(\prod_{k=1}^{n}a_k=a_1\cdots a_n\)。两式中 \(k\) 均依次取 \(1\) 至 \(n\) 的整数。

\section{函数与线性运算}

\(f:A\to B\)表示映射的定义域为\(A\)、陪域为\(B\)；其像为\(f(A)=\{f(x):x\in A\}\)。\((g\circ f)(x)=g(f(x))\)表示先使用\(f\)，再使用\(g\)。\(f^{-1}\)表示已经存在的反函数，与逐点取倒数\(1/f\)分别理解。

线性映射\(T:V\to W\)的核与像分别为
\begin{equation*}
 \ker T=\{v\in V:Tv=0\},\qquad
 \operatorname{im}T=\{Tv:v\in V\}.
\end{equation*}
\(\operatorname{span}(v_1,\ldots,v_k)\)表示这些向量全部线性组合构成的子空间。内积记为\(\langle u,v\rangle\)，其诱导的长度记为\(\|u\|=\sqrt{\langle u,u\rangle}\)；几何向量的长度在正文中也记为\(|u|\)。

矩阵的转置为\(A^T\)，可逆矩阵的逆为\(A^{-1}\)，单位矩阵为\(I\)。行列式记为\(\det A\)。这些记号的作用对象由当前的向量空间与坐标表示确定。

\section{极限与近似}

\(a_n\to a\)表示数列收敛到\(a\)，\(f(x)\to L\)、\(x\to x_0\)表示相应函数极限。余项\(r(h)=o(h^n)\)表示\(r(h)/h^n\to0\)；\(r(h)=O(h^n)\)表示在所讨论极限附近，\(|r(h)|\le C|h|^n\)可由某个固定常数\(C\)控制。

\(C[a,b]\)、\(C^1[a,b]\)、\(C^2[a,b]\)依次表示连续函数、一阶连续可导函数、二阶连续可导函数组成的空间。连续函数的最大范数为\(\|f\|_\infty=\max_{[a,b]}|f|\)，比较两个函数时由\(\|f-g\|_\infty\)衡量整个区间上的最大误差。
